%==============================================================================
% Appendix B: Shadow Map Format
%==============================================================================
\section{Shadow Map Data Structure}
\label{app:shadow_map}

The shadow map is the central data structure that defines how the \ac{PCH}'s linear address space is mapped to physical flash devices. It is stored in two locations for redundancy:
\begin{enumerate}
    \item \textbf{Extension Flash:} At fixed offset \texttt{0x00EE0000} (4 KB region, holds up to 256 entries).
    \item \textbf{24C02 EEPROM:} Full 256-byte map (16 entries $\times$ 16 bytes) for instant boot-time availability.
\end{enumerate}

\subsection{Shadow Map Header (Extension Flash)}
\begin{lstlisting}[style=vbfcStyle, caption={Shadow Map Header in Extension Flash}, label={lst:shadow_header}]
struct shadow_map_header {
    uint32_t magic;           // 0x53484457 ("SHDW")
    uint16_t version;         // Format version (1)
    uint16_t num_entries;     // Number of valid entries (0-256)
    uint32_t crc32;           // CRC-32 of all entries
    uint64_t timestamp;       // Unix timestamp of last update
    uint8_t  reserved[16];    // Future use
};
\end{lstlisting}

\subsection{Shadow Map Entry (16 bytes)}
As defined in \tabref{tab:shadow_map_entry}, each entry maps a contiguous logical address range to a physical source.

\begin{table}[h]
\centering
\caption{Region Type Encoding}
\label{tab:region_type}
\begin{tabular}{@{}cll@{}}
\toprule
\textbf{Value} & \textbf{Name} & \textbf{Behavior} \\
\midrule
0x00 & ORIGINAL & Passthrough to motherboard BIOS flash \\
0x01 & EXTENSION & Read from extension flash at \texttt{physical\_offset} \\
0x02 & HYBRID   & Read from extension flash + apply patch cache \\
0x03 & SYNTHESIZED & Generate response in firmware (e.g., JEDEC ID, Status Reg) \\
0xFF & INVALID  & Entry unused \\
\bottomrule
\end{tabular}
\end{table}

\subsection{Flags Field}
\begin{itemize}
    \item Bit 0 (0x01): \texttt{PATCH\_ENABLED} — Patch cache active for this region (HYBRID only).
    \item Bit 1 (0x02): \texttt{READ\_ONLY} — Write/erase commands blocked for this region.
    \item Bit 2 (0x04): \texttt{CACHE\_LOCKED} — Patch cache entries for this region pinned.
    \item Bits 3-7: Reserved (zero).
\end{itemize}

\subsection{Example Shadow Map Configuration}
For a typical hidden-menu unlock scenario on a \SI{16}{MB} flash:

\begin{table}[h]
\centering
\caption{Example Shadow Map (Hidden Menu Unlock)}
\label{tab:shadow_example}
\begin{tabular}{@{}cccccc@{}}
\toprule
\textbf{\#} & \textbf{Logical Base} & \textbf{Logical Limit} & \textbf{Phys Offset} & \textbf{Type} & \textbf{Flags} \\
\midrule
0 & 0x000000 & 0x0FFFFF & 0x000000 & ORIGINAL & 0x00 \\
1 & 0x100000 & 0x1FFFFF & 0x010000 & HYBRID   & 0x01 \\
2 & 0x200000 & 0xFFFFFF & 0x020000 & EXTENSION & 0x00 \\
\bottomrule
\end{tabular}
\end{table}

\begin{itemize}
    \item Entry 0: First 1 MB (typically PEI/SEC phase) passthrough to original — critical for Boot Guard IBB.
    \item Entry 1: 1--2 MB (DXE volume with setup menu code) from extension flash with patching enabled.
    \item Entry 2: Remaining 14 MB from extension flash (unpatched).
\end{itemize}

\subsection{Shadow Map Lookup Algorithm}
\begin{lstlisting}[style=vbfcStyle, caption={Shadow Map Lookup (C)}, label={lst:shadow_lookup}]
region_type_t shadow_map_lookup(uint32_t logical_addr,
                                uint32_t *phys_offset_out) {
    for (int i = 0; i < shadow_map_hdr.num_entries; i++) {
        entry_t *e = &shadow_map[i];
        if (logical_addr >= e->logical_base &&
            logical_addr <= e->logical_limit) {
            *phys_offset_out = e->physical_offset + 
                               (logical_addr - e->logical_base);
            return e->region_type;
        }
    }
    return REGION_INVALID; // Fallback: original flash
}
\end{lstlisting}

The linear search is acceptable for $\le$16 entries. For larger maps (future), a binary search on sorted \texttt{logical\_base} or a small radix tree can be used.

\subsection{Atomic Shadow Map Updates}
To prevent partial updates from corrupting the map, updates use a double-buffered scheme:
\begin{enumerate}
    \item Write new map to staging area in extension flash.
    \item Compute CRC32, verify.
    \item Write new header with updated \texttt{timestamp} and \texttt{crc32} (single 4 KB sector erase + program).
    \item Update EEPROM copy via I2C (256 bytes, \SI{\approx 5}{ms}).
    \item On boot, arbiter loads map from EEPROM (fast) and verifies against flash copy.
\end{enumerate}